%% ma: L10-B14-0013
%% lop: 10
%% chuong: 5
%% bai: 14
%% dang: 10.14.2
%% loai: TN4
%% muc: TH
%% dap_an: "C"
%% nguon: "chuyen-de-cac-so-dac-trung-cua-mau-so-lieu-khong-ghep-nhom-toan-10-chuong-trinh-moi.pdf (GV. Phan Nhật Linh) (P04), Bài 03 Dạng 1 Phần I Câu 3"
%% kien_thuc: "Lớp 10 Bài 13 (số trung bình, tứ phân vị), Bài 14 (khoảng biến thiên, khoảng tứ phân vị, giá trị bất thường)"
%% trang_thai: chinh-thuc
%% ghi_chu: "Tự duyệt theo yêu cầu giáo viên 22/10/2026."
\begin{ex}
Mẫu số liệu sau đây cho biết số bài hát ở mỗi album trong bộ sưu tập của Bình:
\begin{center}
\begin{tabular}{|c|c|c|c|c|c|c|c|c|c|c|c|c|c|c|c|}\hline
$12$ & $7$ & $10$ & $9$ & $12$ & $9$ & $7$ & $11$ & $10$ & $14$ & $8$ & $6$ & $13$ & $8$ & $11$ & $8$\\\hline
\end{tabular}
\end{center}
Khoảng biến thiên và khoảng tứ phân vị của mẫu số liệu lần lượt là
\choice{$R=8$ và $\Delta_Q=4$}{$R=10$ và $\Delta_Q=3{,}5$}{\True $R=8$ và $\Delta_Q=3{,}5$}{$R=10$ và $\Delta_Q=4$}
\loigiai{Sắp xếp: $6$; $7$; $7$; $8$; $8$; $8$; $9$; $9$; $10$; $10$; $11$; $11$; $12$; $12$; $13$; $14$. $R=14-6=8$.

$Q_1=\dfrac{8+8}2=8$; $Q_3=\dfrac{11+12}2=11{,}5$; $\Delta_Q=3{,}5$. Chọn C.}
\end{ex}
